\section{The polynomial horizon}
\label{sec:T_polynomial}

We now combine the lemmas of
\Cref{sec:anchor_decomposition,sec:anchor_accuracy,sec:anchor_count,sec:anchor_mass}
into a single high-probability bound on $\norm{x_T - V^*(Q)}$. The result
isolates the score-margin condition from the time-horizon: the deviation
splits into a $T$-independent leakage term controlled by $\Delta$ and a
high-probability event controlled by $T$.

\begin{lemma}[Polynomial horizon to reach the decoding margin]\label{lem:T_polynomial}
Fix $Q \in F$ and a confidence parameter $\delta \in (0, 1)$. Suppose
\Cref{ass:anchor_set_accuracy,ass:anchor_emission_prob,ass:score_margin,ass:bounded_value_norms,ass:decoding_existence}
hold and the score-margin constant $\Delta$ of \Cref{ass:score_margin}
satisfies
\begin{align}\label{eq:Delta_condition}
   \Delta \;\ge\; \log\!\frac{4(M + \norm{V^*(Q)})}{p_0\, \gamma(Q)}.
\end{align}
Then for every horizon
\begin{align}\label{eq:T_polynomial_def}
   T \;\ge\; T(Q, \delta)
   \;\coloneqq\; \frac{8}{p_0} \log\!\frac{2}{\delta},
\end{align}
we have
\begin{align}\label{eq:T_polynomial_conclusion}
   \Pr\!\Big[\, \norm{x_T - V^*(Q)} \;\le\; \tfrac{\gamma(Q)}{2}
              \;+\; \varepsilon_{\mathrm{anc}} \,\Big]
   \;\ge\; 1 - \tfrac{\delta}{2}.
\end{align}
\end{lemma}

\begin{proof}
We combine \Cref{lem:anchor_decomposition,lem:anchor_accuracy_bound,lem:anchor_count_lb,lem:anchor_mass_lb}.
The randomness enters only through \Cref{lem:anchor_count_lb}
(anchor-count concentration); conditional on the event
$\Ecal_1 \coloneqq \{|\mathcal A^{\mathrm{traj}}_T| \ge p_0 T / 2\}$,
the remaining algebra is deterministic.

\noindent\textbf{Step 1: anchor-count event.}
By \Cref{lem:anchor_count_lb} applied at confidence $\delta_1 = \delta/2$
with hypothesis $T \ge 8 \log(2/\delta)/p_0$ (which is
\Cref{eq:T_polynomial_def}), $\Pr[\Ecal_1] \ge 1 - \delta/2$.

\noindent\textbf{Step 2: combine anchor decomposition with the accuracy
and mass bounds.}
By \Cref{lem:anchor_decomposition} with the choice
$\mathcal A = \mathcal A^{\mathrm{traj}}_T$,
\begin{align}\label{eq:T_poly_decomp_use}
   \norm{x_T - V^*}
   \;\le\;
   \sum_{k \in \mathcal A^{\mathrm{traj}}_T} w_{T,k}\, \norm{V_k - V^*}
   \;+\;
   \Big(1 - \sum_{k \in \mathcal A^{\mathrm{traj}}_T} w_{T,k}\Big) \cdot D_T.
\end{align}
By \Cref{lem:anchor_accuracy_bound}, the first summand is bounded by
\begin{align*}
   \sum_{k \in \mathcal A^{\mathrm{traj}}_T} w_{T,k}\, \norm{V_k - V^*}
   \;\le\;
   \varepsilon_{\mathrm{anc}} \cdot
   \sum_{k \in \mathcal A^{\mathrm{traj}}_T} w_{T,k}
   \;\le\;
   \varepsilon_{\mathrm{anc}}.
\end{align*}
By the triangle inequality and \Cref{ass:bounded_value_norms},
$D_T = \max_k \norm{V_k - V^*} \le \max_k \norm{V_k} + \norm{V^*} \le M + \norm{V^*}$.

By \Cref{lem:anchor_mass_lb} applied on $\Ecal_1$,
\begin{align*}
   1 - \sum_{k \in \mathcal A^{\mathrm{traj}}_T} w_{T,k}
   \;\le\;
   \frac{|\{1,\dots,T\} \setminus \mathcal A^{\mathrm{traj}}_T|}
        {|\mathcal A^{\mathrm{traj}}_T|}
   \cdot e^{-\Delta}
   \;\le\;
   \frac{T}{p_0 T / 2} \cdot e^{-\Delta}
   \;=\;
   \frac{2}{p_0}\, e^{-\Delta},
\end{align*}
where the last inequality uses
$|\{1,\dots,T\}\setminus \mathcal A^{\mathrm{traj}}_T| \le T$ and
$|\mathcal A^{\mathrm{traj}}_T| \ge p_0 T / 2$ (the latter is $\Ecal_1$).

Combining the three pieces, on $\Ecal_1$,
\begin{align}\label{eq:T_poly_combined}
   \norm{x_T - V^*}
   \;\le\;
   \varepsilon_{\mathrm{anc}}
   \;+\;
   \frac{2}{p_0}\, e^{-\Delta} \cdot (M + \norm{V^*}).
\end{align}

\noindent\textbf{Step 3: apply the score-margin hypothesis.}
The hypothesis \Cref{eq:Delta_condition} is equivalent to
$\tfrac{2}{p_0} e^{-\Delta}(M + \norm{V^*}) \le \gamma(Q)/2$.
Indeed,
$e^{-\Delta} \le p_0 \gamma(Q) / (4(M + \norm{V^*}))$ rearranges to
$\Delta \ge \log(4(M + \norm{V^*})/(p_0 \gamma(Q)))$; multiplying by
$\tfrac{2}{p_0}(M + \norm{V^*})$ gives
$\tfrac{2}{p_0} e^{-\Delta}(M + \norm{V^*}) \le \gamma(Q)/2$.
Substituting into \Cref{eq:T_poly_combined} yields, on $\Ecal_1$,
\begin{align*}
   \norm{x_T - V^*}
   \;\le\; \varepsilon_{\mathrm{anc}} + \tfrac{\gamma(Q)}{2}.
\end{align*}

\noindent\textbf{Step 4: probability budget.}
Combining Step 1 ($\Pr[\Ecal_1] \ge 1 - \delta/2$) with the
deterministic bound on $\Ecal_1$ from Step 3,
$\Pr[\norm{x_T - V^*} \le \gamma(Q)/2 + \varepsilon_{\mathrm{anc}}]
\ge \Pr[\Ecal_1] \ge 1 - \delta/2$, which is
\Cref{eq:T_polynomial_conclusion}. This completes the proof.
\end{proof}

\begin{remark}\label{rem:T_polynomial_remark}
The lemma cleanly separates the two roles of the assumptions: the
score-margin condition \Cref{eq:Delta_condition} is a $T$-independent
property of the trained model, while the polynomial horizon
\Cref{eq:T_polynomial_def} controls the concentration failure budget.
If the trained model's $\Delta$ does not satisfy
\Cref{eq:Delta_condition}, no horizon $T$ can drive the deviation below
$\gamma(Q)/2 + \varepsilon_{\mathrm{anc}}$ in this proof --- the leakage
term $(2/p_0) e^{-\Delta}(M + \norm{V^*})$ does not shrink with $T$. The
score-margin condition is the load-bearing ``trained-model''
requirement; the horizon $T$ only suppresses the multiplicative-Chernoff
concentration tail.
\end{remark}
